In this chapter, we study techniques that further exploit the geometry
of convex functions and associated norms. Many of these techniques
are effective in practice for large scale problems. We recall the
following relevant definitions.

\section{Norms and Local Metrics}

Any centrally symmetric convex body $K$ induces the following norm:
\[
\norm x_{K}=\inf\left\{ t:x\in tK\right\} .
\]
For the usual Euclidean norm, the convex body is the Euclidean ball,
$B_{2}^{n}$. Similarly, for an $\ell_{p}$ norm, the convex body
is the unit $\ell_{p}$ ball. 

It is often useful to consider affine transformations, and the (semi)norms
they induce, e.g., for a PSD matrix $A$ (would be an induced norm
in the case when A is PD), we can define the associated norm as follows:
\[
\norm x_{A}=\sqrt{x^{\top}Ax}.
\]
The convex body (unit ball) of this norm is an ellipsoid centered
at zero and defined by the matrix $A$. 

As we have encountered previously in this book, convex sets and functions
have duals. We recall the dual (polar) of a convex set $K$: 
\[
K^{*}=\left\{ y:\forall x\in K,\left\langle x,y\right\rangle \le1\right\} .
\]

The dual norm for $\norm ._{K}$ is $\norm ._{K^{*}}$. We can state
a generalized Cauchy-Schwarz inequality.
\begin{fact}
For $x,y\in\R^{n}$ and any centrally symmetric convex body $K$,
we have $\left\langle x,y\right\rangle \le\norm x_{K}\norm y_{K^{*}}.$
\end{fact}

In this chapter, an important idea will be \emph{local} norms, i.e.,
at each point $x$ in the domain, there could be a different norm.
Indeed, in a Riemannian metric ${\cal M}$, for every $x\in{\cal M},$
there is a matrix $A(x)$ s.t. the norm at $x$ is defined as $\norm v_{x}=\sqrt{v^{\top}A(x)v}.$
A special class of Riemannian metrics of particular interest for us
will be \emph{Hessian }metrics (corresponding to Hessian manifolds).
Here the matrix defining the local norm is the Hessian of a convex,
twice-differentiable function, i.e., 
\[
A(x)=\nabla^{2}\phi(x)
\]
for some convex function $\phi$. 


\section{Mirror Descent\label{sec:mirror_descent}}

The cutting plane method is well-suited for minimizing non-smooth
convex functions with high accuracy. However, its relatively large
polynomial time complexity and quadratic space requirement are not
favorable for large scale problems. Here we discuss a different approach
to minimize a non-smooth function with low accuracy.

\subsection{Subgradient method}

In this section, we consider the constrained non-smooth minimization
problem $\min_{x\in\mathcal{D}}f(x)$. Recall how to define gradient
for non-smooth functions:
\begin{defn}
For any convex function $f$, we define $\partial f(x)$ be the set
of vectors $g$ such that
\[
f(y)\geq f(x)+g^{\top}(y-x)\text{ for all }y\in\Rn.
\]
Such a vector $g$ is called a \emph{subgradient} of $f$ at $x$.
\end{defn}

\begin{algorithm2e}[H]

\caption{$\mathtt{SubgradientMethod}$}

\SetAlgoLined

\textbf{Input:} Initial point $x^{(0)}\in\Rn$, step size $h>0$.

\For{$k=0,1,\cdots,T-2$}{

Pick any $g^{(k)}\in\partial f(x^{(k)})$.

$y^{(k+1)}\leftarrow x^{(k)}-h\cdot g^{(k)}$.

$x^{(k+1)}\leftarrow\pi_{\mathcal{D}}(y^{(k+1)})$ where $\pi_{\mathcal{D}}(y)=\arg\min_{x\in\mathcal{D}}\|x-y\|_{2}$

}

\Return $\frac{1}{T}\sum_{k=0}^{T-1}x^{(k)}$.

\end{algorithm2e}

Note that we return the average of all iterates, rather than the last
iterate, as the average is better behaved in the worst case. To analyze
the algorithm, we first need the following Pythagorean theorem. This
shows that when we project a point to a convex set, we get closer
to each point in the convex set, and how much closer depends on how
much we move.
\begin{lem}[Pythagorean Theorem]
\label{lem:Pythagorean}Given a convex set $\mathcal{D}$ and a point
$y$, let $\pi(y)=arg\min_{x\in\mathcal{D}}\|x-y\|_{2}$. For any
$z\in\mathcal{D}$, we have that
\[
\|z-\pi(y)\|_{2}^{2}+\|\pi(y)-y\|_{2}^{2}\leq\|z-y\|_{2}^{2}.
\]
\end{lem}

\begin{proof}
Let $h(t)=\|(\pi(y)+t(z-\pi(y)))-y\|^{2}$. Since $\pi(y)$ is the
closest point to $y$, $h(t)$ must be minimized at $t=0$. So, we
have that $h'(0)\geq0$. i.e.,

\[
h'(0)=(\pi(y)-y)^{\top}(z-\pi(y))\geq0.
\]

(If $y$ is in the set, then the statement is trivial since $\pi(y)=y$.
Otherwise, since $\pi(y)$ is the closest point in ${\cal D}$ to
$y$, the hyperplane normal to $\pi(y)-y$ through $\pi(y)$ separates
$y$ from $z$.). Hence,
\begin{align*}
\|z-y\|_{2}^{2} & =\|z-\pi(y)+\pi(y)-y\|_{2}^{2}\\
 & =\|z-\pi(y)\|_{2}^{2}+2(z-\pi(y))^{\top}(\pi(y)-y)+\|\pi(y)-y\|_{2}^{2}\\
 & \geq\|z-\pi(y)\|_{2}^{2}+\|\pi(y)-y\|_{2}^{2}.
\end{align*}
\end{proof}
Now, we are ready to analyze the subgradient method. It basically
involves tracking the squared distance to the optimum, $\|x^{(k+1)}-x^{*}\|_{2}^{2}$.
\begin{thm}
\label{thm:projected_gradient_descent}Let $f$ be a convex function
that is $G$-Lipschitz in $\ell_{2}$ norm. After $T$ steps, the
subgradient method outputs a point $x$ such that
\[
f(x)\leq f(x^{*})+\frac{\|x^{(0)}-x^{*}\|_{2}^{2}}{2hT}+\frac{h}{2}G^{2}
\]
where $x^{*}$ is any point that minimizes $f$ over ${\cal D}$.
\end{thm}

\begin{rem*}
If the distance $\|x^{(0)}-x^{*}\|$ and the Lipschitz constant $G$
are known, we can pick $h=\frac{\|x^{(0)}-x^{*}\|_{2}}{G\sqrt{T}}$
and get
\[
f(x)\leq f(x^{*})+\frac{G\cdot\|x^{(0)}-x^{*}\|_{2}}{\sqrt{T}}.
\]
\end{rem*}
\begin{proof}
Let $x^{*}$ be any point that minimizes $f$. Then, we have
\begin{align*}
\|x^{(k+1)}-x^{*}\|_{2}^{2} & =\|\pi_{\mathcal{D}}(y^{(k+1)})-x^{*}\|_{2}^{2}\\
 & \leq\|y^{(k+1)}-x^{*}\|_{2}^{2}\\
 & =\|x^{(k)}-hg^{(k)}-x^{*}\|_{2}^{2}\\
 & =\|x^{(k)}-x^{*}\|_{2}^{2}-2h\left\langle g^{(k)},x^{(k)}-x^{*}\right\rangle +h^{2}\|g^{(k)}\|_{2}^{2}.
\end{align*}
where we used Lemma \ref{lem:Pythagorean} in the inequality. Since
$x^{*}$ lies on the $-g^{(k)}$ direction, we expect $x^{(k+1)}$
is closer to $x^{*}$ than $x^{(k)}$ if the step size is small enough
(or if we ignore the second order term $\|g^{(k)}\|_{2}^{2}$). To
bound the distance improvement, we apply the definition of subgradient
and get
\[
f(x^{*})\geq f(x^{(k)})+\left\langle g^{(k)},x^{*}-x^{(k)}\right\rangle .
\]
Therefore, we have that
\begin{align*}
\|x^{(k+1)}-x^{*}\|_{2}^{2} & \leq\|x^{(k)}-x^{*}\|_{2}^{2}-2h\cdot(f(x^{(k)})-f(x^{*}))+h^{2}\|g^{(k)}\|_{2}^{2}\\
 & \leq\|x^{(k)}-x^{*}\|_{2}^{2}-2h\cdot(f(x^{(k)})-f(x^{*}))+h^{2}G^{2}.
\end{align*}
Note that this equation shows that if the error $f(x^{(k)})-f(x^{*})$
is larger, then we move faster towards the optimum. Rearranging the
terms, we have
\[
f(x^{(k)})-f(x^{*})\leq\frac{1}{2h}\left(\|x^{(k)}-x^{*}\|_{2}^{2}-\|x^{(k+1)}-x^{*}\|_{2}^{2}\right)+\frac{h}{2}G^{2}.
\]
We sum over all iterations, to get
\begin{align*}
\frac{1}{T}\sum_{k=0}^{T-1}(f(x^{(k)})-f(x^{*})) & \leq\frac{1}{T}\cdot\frac{1}{2h}(\|x^{(0)}-x^{*}\|_{2}^{2}-\|x^{(T)}-x^{*}\|_{2}^{2})+\frac{h}{2}G^{2}\\
 & \leq\frac{\|x^{(0)}-x^{*}\|_{2}^{2}}{2hT}+\frac{h}{2}G^{2}.
\end{align*}
The result follows from observing that for a convex function, 
\[
f\left(\frac{1}{T}\sum_{k=0}^{T-1}x^{(k)}\right)-f(x^{*})\leq\frac{1}{T}\sum_{k=0}^{T-1}(f(x^{(k)})-f(x^{*})).
\]
\end{proof}


\subsection{Intuition of Mirror Descent}

Consider using the subgradient method above to minimize $f(x)=\sum_{i}|x_{i}|$
over the unit ball $B(0,1)$. The subgradient of $f$ is given by
$\text{sign}(x)$ (assuming $x_{i}\neq0$ for all $i$). Therefore,
we have that the Lipschitz constant is bounded by the Euclidean norm
of a vector with $\pm1$ entries, i.e., $G=\sqrt{n}$ and the set
is contained in a ball of radius $R=1$. Hence, by the main theorem
of the previous section, the error is bounded by $\sqrt{\frac{n}{T}}$
after $T$ steps, which grows with the dimension. Intuitively, the
dimension dependence comes from the fact that we must take a tiny
step size to avoid changing the variables too much. For example, if
we take a constant step size $h$ and start at the point $x=(1,\frac{1}{n},\frac{1}{n},\cdots,\frac{1}{n})$,
then we will get a point with $x_{i}$ constant in all directions
$i\neq1$, and this increases the function $f$ dramatically from
$\Theta(1)$ to $\Theta(nh)$. Therefore, to get constant error, we
need to take step size $h\le\frac{1}{n}$ and hence we need $\Omega(n)$
iterations.

Conceptually, the step $x=x-\eta g$ does not make sense either. Imagine
the problem is infinitely dimensional. Note that $f(x)<+\infty$ for
any 
\[
x\in\ell_{1}=\{x\in\R^{\mathbb{N}}:\sum_{i}|x_{i}|<\infty\}.
\]
On the other hand, its gradient $g$ lives in $\ell_{\infty}$ space;
the dual space of $\ell_{1}$ is $\ell_{\infty}$ (in general, $\ell_{p}$
is dual to $\ell_{q}$ where $(1/p)+(1/q)=1$). Since $x$ and $g$
are not in the same space, the term $x-\eta g$ does not make sense.
More precisely, we have the following tautology (directly follows
from the definition of dual space, namely the set of all linear maps
in the original space).
\begin{defn}
A \emph{Banach space }over the reals is a vector space over the reals
together with a norm that defines a complete metric space, i.e., for
any Cauchy sequence, $X=(x_{i})_{i=1}^{\infty}$, there is a vector
$x$ s.t. $\lim_{n\rightarrow\infty}\norm{x_{n}-x}=0$. 
\end{defn}

\begin{fact}
\label{lem:grad_dual}Given any Banach space $\mathcal{D}$ over the
reals and a continuously differentiable function $f$ from $\mathcal{D}$
to $\R$, its gradient $\nabla f(x)\in\mathcal{D}^{*}$ for any $x$.
\end{fact}

In general, if the function $f$ is on the primal space $\mathcal{D}$,
then the gradient $g$ lives in the dual space $\mathcal{D}^{*}$.
Therefore, we need to map $x$ from the primal space $\mathcal{D}$
to the dual space $\mathcal{D}^{*}$, update its position, then map
the point back to the original space $\mathcal{D}$.

In fact, Lemma \ref{lem:grad_dual} gives us one such map, $\nabla f$.
Consider the following algorithm: Starting at $x$. We use $\nabla f(x)$
to map $x$ to the dual space $y=\nabla f(x)$. Then, we apply the
gradient step on the dual $y^{\new}=y-\nabla f(x)$ and map it back
to the primal space, namely finding $x^{\new}$ such that $\nabla f(x^{\new})=y^{\new}$.
Note that $y^{\new}=0$ and hence $x^{\new}$ is exactly a minimizer
of $f$. So, if we can map it back, this algorithm solves the problem
in one step. Unfortunately, the task of mapping it back is exactly
our original problem.

Instead of using the same $f$, mirror descent uses some other convex
function $\Phi$, called the \emph{mirror map}. For constrained problems,
the mirror map may not bring the point back to a point in $\mathcal{D}$.
Naively, one may consider the algorithm 
\begin{align*}
\nabla\Phi(y^{(t+1)}) & =\nabla\Phi(x^{(t)})-h\cdot\nabla f(x^{(t)}),\\
x^{(t+1)} & =\arg\min_{x\in\mathcal{D}}\|x-y^{(t+1)}\|_{2}.
\end{align*}
Note that the first step of finding $y^{(t+1)}$ involves solving
an optimization problem. We will show how to do this optimization
later (see \ref{eq:mirror_descent}) but with a proper formulation
of the algorithm which takes into account the distance as measured
by the mirror map $\Phi$.
\begin{defn}
For any strictly convex function $\Phi$, we define the Bregman divergence
as
\[
D_{\Phi}(y,x)=\Phi(y)-\Phi(x)-\left\langle \nabla\Phi(x),y-x\right\rangle .
\]

Note that $D_{\Phi}(y,x)$ is the error of the first order Taylor
expansion of $\Phi$ at $x$. Due to the convexity of $\Phi$, we
have that $D_{\Phi}(y,x)\geq0$. Also, we note that $D_{\Phi}(y,x)$
is convex in $y$, but not necessarily in $x$.
\end{defn}

\begin{example}
$D_{\Phi}(y,x)=\|y-x\|^{2}$ for $\Phi(x)=\|x\|^{2}$. $D_{\Phi}(y,x)=\sum_{i}y_{i}\log\frac{y_{i}}{x_{i}}-\sum y_{i}+\sum x_{i}$
for $\Phi(x)=\sum x_{i}\log x_{i}$.
\end{example}

\begin{algorithm2e}[H]

\caption{$\mathtt{MirrorDescent}$}

\SetAlgoLined

\textbf{Input:} Initial point $x^{(0)}=\arg\min_{x\in\mathcal{D}}\Phi(x)$,
step size $h>0$.

\For{$k=0,1,\cdots,T-2$}{

Pick any $g^{(k)}\in\partial f(x^{(k)})$.

\tcp{As shown in (\ref{eq:mirror_descent}), the next 2 steps can be implemented by an optimization over $D_{\Phi}$.}

Find $y^{(k+1)}$ such that $\nabla\Phi(y^{(k+1)})=\nabla\Phi(x^{(k)})-h\cdot g^{(k)}$.

$x^{(k+1)}\in\pi_{\mathcal{D}}^{\Phi}(y^{(k+1)})$ where $\pi_{\mathcal{D}}^{\Phi}(y)=\arg\min_{x\in\mathcal{D}}D_{\Phi}(x,y)$.

}

\Return $\frac{1}{T}\sum_{k=0}^{T-1}x^{(k)}$.

\end{algorithm2e}

The Mirror Descent step can be written as follows. In the second step
below, we use the fact that the optimization is only over $x$ (and
hence all terms in $y$ can be ignored):
\begin{align}
x^{(k+1)} & =\arg\min_{x\in\mathcal{D}}D_{\Phi}(x,y^{(k+1)})\nonumber \\
 & =\arg\min_{x\in\mathcal{D}}\Phi(x)-\Phi(y^{(k+1)})-\nabla\Phi(y^{(k+1)})^{\top}(x-y^{(k+1)})\nonumber \\
 & =\arg\min_{x\in\mathcal{D}}\Phi(x)-\nabla\Phi(y^{(k+1)})^{\top}x\nonumber \\
 & =\arg\min_{x\in\mathcal{D}}\Phi(x)-\nabla\Phi(x^{(k)})^{\top}x+hg^{(k)\top}x\nonumber \\
 & =\arg\min_{x\in\mathcal{D}}hg^{(k)\top}x+D_{\Phi}(x,x^{(k)}).\label{eq:mirror_descent}
\end{align}
Note that this is a natural generalization of $x^{(k+1)}=\arg\min_{x\in\mathcal{D}}hg^{(k)\top}x+\|x-x^{(k)}\|^{2}.$

\subsection{Analysis of Mirror Descent}
\begin{lem}[Pythagorean Theorem]
\label{lem:Pythagorean2}Given a convex set $\mathcal{D}$ and a
point $y$, let $\pi(y)=\arg\min D_{\Phi}(x,y)$. For any $z\in\mathcal{D}$,
we have that
\[
D_{\Phi}(z,\pi(y))+D_{\Phi}(\pi(y),y)\leq D_{\Phi}(z,y).
\]
\end{lem}

\begin{proof}
Let $h(t)=D_{\Phi}(\pi(y)+t(z-\pi(y)),y)$. Since $h(t)$ is minimized
at $t=0$, we have that $h'(0)\geq0$. Hence, we have
\[
h'(0)=(\nabla\Phi(\pi(y))-\nabla\Phi(y))^{\top}(z-\pi(y))\geq0.
\]
Hence,
\begin{align*}
D_{\Phi}(z,y) & =D_{\Phi}(z,\pi(y))+2(z-\pi(y))^{\top}(\nabla\Phi(\pi(y))-\nabla\Phi(y))+D_{\Phi}(\pi(y),y)\\
 & \geq D_{\Phi}(z,\pi(y))+D_{\Phi}(\pi(y),y).
\end{align*}
\end{proof}
\begin{thm}
\label{thm:mirror_descent}Let $f$ be a $G$-Lipschitz convex function
on ${\cal D}$ with respect to some norm $\|\cdot\|$. Let $\Phi$
be a $\rho$-strongly convex function on $\mathcal{D}$ with respect
to $\|\cdot\|$ with squared diameter $R^{2}=\sup_{x\in\mathcal{D}}\Phi(x)-\Phi(x^{(0)})$.
Then, mirror descent outputs $x$ such that
\[
f(x)-\min_{x}f(x)\leq\frac{R^{2}}{hT}+\frac{h}{2\rho}G^{2}.
\]
\end{thm}

\begin{rem}
We say a function $f$ is $\rho$ strongly convex with respect to
the norm $\|\cdot\|$ if for any $x,y$, we have 
\[
f(y)\geq f(x)+\left\langle \nabla f(x),y-x\right\rangle +\frac{\rho}{2}\|y-x\|^{2}.
\]
The usual strong convexity is with respect to the Euclidean norm.
\end{rem}

\begin{rem}
Picking $h=\frac{R}{G}\sqrt{\frac{2\rho}{T}}$, we get the rate $f(x)\leq\min_{x}f(x)+GR\sqrt{\frac{2}{\rho T}}.$
\end{rem}

\begin{proof}
This proof is a complete ``mirror'' of the proof in Theorem \ref{thm:projected_gradient_descent}.
\begin{align*}
D_{\Phi}(x^{*},x^{(k+1)}) & \leq D_{\Phi}(x^{*},y^{(k+1)})-D_{\Phi}(x^{(k+1)},y^{(k+1)})\\
 & =D_{\Phi}(x^{*},x^{(k)})+(x^{*}-x^{(k)})^{\top}(\nabla\Phi(x^{(k)})-\nabla\Phi(y^{(k+1)}))+D_{\Phi}(x^{(k)},y^{(k+1)})-D_{\Phi}(x^{(k+1)},y^{(k+1)})\\
 & =D_{\Phi}(x^{*},x^{(k)})-h\cdot\left\langle g^{(k)},x^{(k)}-x^{*}\right\rangle +D_{\Phi}(x^{(k)},y^{(k+1)})-D_{\Phi}(x^{(k+1)},y^{(k+1)})
\end{align*}
where we used Lemma \ref{lem:Pythagorean2} in the inequality. Using
the definition of subgradient, we have that
\[
f(x^{*})\geq f(x^{(k)})+\left\langle g^{(k)},x^{*}-x^{(k)}\right\rangle .
\]
Therefore, we have that
\begin{align*}
D_{\Phi}(x^{*},x^{(k+1)}) & \leq D_{\Phi}(x^{*},x^{(k)})-h\cdot(f(x^{(k)})-f(x^{*}))+D_{\Phi}(x^{(k)},y^{(k+1)})-D_{\Phi}(x^{(k+1)},y^{(k+1)}).
\end{align*}
Next we note that
\begin{align*}
 & D_{\Phi}(x^{(k)},y^{(k+1)})-D_{\Phi}(x^{(k+1)},y^{(k+1)})\\
= & \Phi(x^{(k)})-\Phi(x^{(k+1)})-\nabla\Phi(y^{(k+1)})^{\top}(x^{(k)}-x^{(k+1)})\\
\leq & \left\langle \nabla\Phi(x^{(k)})-\nabla\Phi(y^{(k+1)}),x^{(k)}-x^{(k+1)}\right\rangle -\frac{\rho}{2}\|x^{(k)}-x^{(k+1)}\|^{2}\\
= & h\cdot\left\langle g^{(k)},x^{(k)}-x^{(k+1)}\right\rangle -\frac{\rho}{2}\|x^{(k)}-x^{(k+1)}\|^{2}\\
\leq & hG\|x^{(k)}-x^{(k+1)}\|-\frac{\rho}{2}\|x^{(k)}-x^{(k+1)}\|^{2}\\
\leq & \frac{(hG)^{2}}{2\rho}.
\end{align*}
Hence, we have
\begin{align*}
D_{\Phi}(x^{*},x^{(k+1)}) & \leq D_{\Phi}(x^{*},x^{(k)})-h\cdot(f(x^{(k)})-f(x^{*}))+\frac{(hG)^{2}}{2\rho}.
\end{align*}

Rearranging the terms, we have
\[
f(x^{(k)})-f(x^{*})\leq\frac{1}{h}\left(D_{\Phi}(x^{*},x^{(k)})-D_{\Phi}(x^{*},x^{(k+1)})\right)+\frac{hG^{2}}{2\rho}.
\]
Summing over all iterations, we have
\begin{align*}
\frac{1}{T}\sum_{k=0}^{T-1}(f(x^{(k)})-f(x^{*})) & \leq\frac{1}{T}\cdot\frac{1}{h}(D_{\Phi}(x^{*},x^{(0)})-D_{\Phi}(x^{*},x^{(T)}))+\frac{h}{2\rho}G^{2}\\
 & \leq\frac{1}{hT}D_{\Phi}(x^{*},x^{(0)})+\frac{h}{2\rho}G^{2}\\
 & \leq\frac{R^{2}}{hT}+\frac{h}{2\rho}G^{2}.
\end{align*}
Using the fact that $x^{(0)}$ was chosen as a minimum of $\Phi(x)$,
\begin{align*}
D_{\Phi}(x^{*},x^{(0)})= & (\Phi(x^{*})-\Phi(x^{(0)})-\left\langle \nabla\Phi(x^{(0)}),x^{*}-x^{(0)}\right\rangle )\\
= & \Phi(x^{*})-\Phi(x^{(0)})\\
\le & R^{2}
\end{align*}

The result follows from the convexity of $f$.
\end{proof}

\subsection{Multiplicative Weight Update}

In this section, we discuss mirror descent under the map $\Phi(x)=\sum x_{i}\log x_{i}$
with the convex set being the simplex $\mathcal{D}=\{x_{i}\geq0,\sum x_{i}=1\}$.

\paragraph{Step Formula}

As we showed in (\ref{eq:mirror_descent}), we have that
\[
x^{(k+1)}=\arg\min_{x\in\mathcal{D}}hg^{(k)\top}x+D_{\Phi}(x,x^{(k)}).
\]
Note that
\begin{align*}
D_{\Phi}(x,x^{(k)}) & =\sum x_{i}\log x_{i}-\sum x_{i}^{(k)}\log x_{i}^{(k)}-\sum_{i}(1+\log x_{i}^{(k)})(x_{i}-x_{i}^{(k)})\\
 & =\sum x_{i}\log\frac{x_{i}}{x_{i}^{(k)}}
\end{align*}
where we used that $\sum_{i}x_{i}=\sum_{i}x_{i}^{(k)}$. Hence, the
step is simply 
\[
x^{(k+1)}=\arg\min_{\sum x_{i}=1,x_{i}\geq0}hg^{(k)\top}x+\sum x_{i}\log\frac{x_{i}}{x_{i}^{(k)}}.
\]
Note that the optimality condition is given by
\[
hg_{i}^{(k)}+\log\frac{x_{i}^{(k+1)}}{x_{i}^{(k)}}+1-\lambda=0
\]
for some Lagrangian multiplier $\lambda$. Rewriting, we have
\[
x_{i}^{(k+1)}=e^{-hg_{i}^{(k)}}x_{i}^{(k)}/Z
\]
for some normalization constant $Z$. Note that this algorithm multiplies
the current $x$ with a multiplicative factor and hence it is also
called multiplicative weight update.

\paragraph{Strong Convexity}

To bound the strong convexity parameter, we note that
\[
\Phi(y)-\Phi(x)-\left\langle \nabla\Phi(x),y-x\right\rangle =\frac{1}{2}(y-x)^{\top}\nabla^{2}\Phi(\zeta)(y-x)
\]
for some $\zeta$ between $x$ and $y$. Since $\nabla^{2}\Phi(\zeta)=\frac{1}{\zeta}$,
we have
\[
(y-x)^{\top}\nabla^{2}\Phi(\zeta)(y-x)=\sum\frac{(y_{i}-x_{i})^{2}}{\zeta_{i}}\geq\frac{(\sum_{i}|y_{i}-x_{i}|)^{2}}{\sum_{i}\zeta_{i}}=(\sum_{i}|y_{i}-x_{i}|)^{2}
\]
where we used that $\sum_{i}x_{i}=\sum_{i}y_{i}=1$ and so $\sum_{i}\zeta_{i}=1$.
Hence, 
\[
\Phi(y)-\Phi(x)-\left\langle \nabla\Phi(x),y-x\right\rangle \geq\frac{1}{2}\|y-x\|_{1}^{2}.
\]
Therefore, $\Phi$ is $1$-strongly convex in $\|\cdot\|_{1}$. Hence,
$\rho=1$.

\paragraph{Diameter}

Direct calculation shows that $-\log n\leq\Phi(x)\leq0$. We start
at $x^{(0)}=\frac{1}{n}(1,\ldots,1)^{T}$. Hence, $R^{2}=\log n$.

\paragraph{Result}
\begin{thm}
Let $f$ be a $1$-Lipschitz function on $\|\cdot\|_{1}$. Then, mirror
descent with the mirror map $\Phi(x)=\sum_{i}x_{i}\log x_{i}$.
\[
f(x^{T})-\min_{x}f(x)\leq\sqrt{\frac{2\log n}{T}}.
\]
\end{thm}

In comparison, projected gradient descent had the bound of $\sqrt{\frac{n}{T}}$.


